\chapter{Speed Estimation}

\section{Algorithmic Formulation}
Speed estimation in \texttt{ai/speed\_estimator.py} operates on tracked object centroids across video frames. The mathematical pipeline comprises four stages:

\subsection{1. Euclidean Pixel Displacement}
Given previous coordinate $(x_1, y_1)$ at monotonic timestamp $t_1$ and current coordinate $(x_2, y_2)$ at timestamp $t_2$:
\begin{align}
\Delta x &= x_2 - x_1, \quad \Delta y = y_2 - y_1 \\
d_{\text{px}} &= \sqrt{(\Delta x)^2 + (\Delta y)^2}
\end{align}

\subsection{2. Heuristic Metric Conversion and Temporal Normalization}
Using a fixed uncalibrated heuristic conversion scale ($K_{\text{scale}} = 0.05\,\text{meters/pixel}$):
\begin{align}
d_{\text{meters}} &= d_{\text{px}} \times K_{\text{scale}} \\
\Delta t &= t_2 - t_1 \quad (\text{seconds}) \\
v_{\text{m/s}} &= \frac{d_{\text{meters}}}{\Delta t} \\
v_{\text{km/h}} &= v_{\text{m/s}} \times 3.6
\end{align}

\section{Robustness Boundaries and Anomaly Protection}
In real video streams, coordinate jitter and track interruptions introduce severe numerical instability:
\begin{itemize}[leftmargin=*]
    \item \textbf{Monotonic Timestamp Guarantee}: $\Delta t$ is computed via \texttt{time.monotonic()} or video frame PTS. If $\Delta t \le 0$, the function returns $0.0\,\text{km/h}$ immediately, preventing division by zero.
    \item \textbf{Tracking Drop Resets}: If a vehicle track disappears or drops for longer than $\text{MAX\_TRACK\_AGE\_SECONDS} = 3.0\,\text{seconds}$, historical coordinates are purged to avoid calculating speed across long pauses.
    \item \textbf{Implausible Glitch Capping}: Single-frame bounding box coordinate jitter or identity swaps can produce momentary artificial spikes. Speeds are hard-capped at $\text{MAX\_PLAUSIBLE\_SPEED\_KMPH} = 180.0\,\text{km/h}$.
    \item \textbf{Memory Leak Protection}: The tracker enforces $\text{MAX\_STORED\_TRACKS} = 250$. When memory thresholds are reached, tracks older than 3.0 seconds are pruned automatically.
\end{itemize}

\section{Why Uncalibrated Pixel Estimation Is Not Real Radar Speed}
\begin{limitationbox}[Camera Perspective and Lens Distortion]
Standard monocular traffic cameras view the roadway through perspective projection: parallel road lane markings converge towards a vanishing point. As a consequence:
\begin{enumerate}
    \item 100 pixels in the foreground near the camera may represent $5\,\text{meters}$ of physical asphalt.
    \item 100 pixels in the background near the horizon may represent $80\,\text{meters}$ of physical asphalt.
\end{enumerate}
Applying a flat scalar ($0.05\,\text{m/px}$) without a 4-point homography calibration matrix $H \in \mathbb{R}^{3 \times 3}$ rectifying the ground plane causes severe speed underestimation in the background and overestimation in the foreground. Therefore, BengaluruTrafficAI explicitly labels all velocity outputs as relative motion estimates rather than legally certified radar measurements.
\end{limitationbox}
